\documentclass{article}
\usepackage[margin=1in]{geometry}
\usepackage{amsmath,amssymb,amsthm}
\usepackage{booktabs,longtable}
\usepackage{hyperref}

\newtheorem{definition}{Definition}
\newtheorem{theorem}[definition]{Theorem}
\newtheorem{lemma}[definition]{Lemma}
\newtheorem{proposition}[definition]{Proposition}
\newtheorem{convention}[definition]{Convention}
\newtheorem{remark}[definition]{Remark}

\newcommand{\source}[1]{\par\noindent\textbf{Source.} #1}
\newcommand{\lean}[1]{\par\noindent\textbf{Lean target.} \texttt{\detokenize{#1}}}
\newcommand{\status}[1]{\par\noindent\textbf{Status.} \texttt{\detokenize{#1}}}
\newcommand{\depends}[1]{\par\noindent\textbf{Depends on.} #1}
\newcommand{\cA}{\mathcal A}
\newcommand{\cH}{\mathcal H}
\newcommand{\cM}{\mathcal M}
\newcommand{\cN}{\mathcal N}
\newcommand{\cB}{\mathcal B}
\newcommand{\OmegaVec}{\Omega}

\begin{document}

\section*{Blueprint 010: Hilbert, GNS, and Represented-Region Layer}

\subsection*{Scope}

This note promotes the third remaining blueprint block from
\path{agents/2026-06-22-remaining-blueprint-suite-goal.md}.  It supports
targets
\[
  \texttt{HILB-001},\quad
  \texttt{GNS-001},\quad
  \texttt{GNS-004},\quad
  \texttt{GNS-002},\quad
  \texttt{GNS-003},\quad
  \texttt{VN-001},\quad
  \texttt{VN-004},\quad
  \texttt{VN-005},\quad
  \texttt{VN-006},\quad
  \texttt{VN-DCT-001},\quad
  \texttt{VN-007},\quad
  \texttt{VN-008},\quad
  \texttt{PROJ-LIM-001},\quad
  \texttt{VN-003}.
\]
It supplies the Hilbert-space and representation layer needed before any
infinite-region von Neumann algebra or factor-type statement is made.  It does
not add Lean declarations, Aristotle packets, Murray--von Neumann comparison,
trace theory, or type-classification predicates.

\subsection*{Takesaki Statement Transcripts}

The local Takesaki PDFs are scanned.  The statement locators below were checked
by rendering and OCR of the cited pages into `/tmp/formalization-ocr`.

\begin{theorem}[GNSCyclicRepresentation, Takesaki I Theorem I.9.14, PDF pp.~59--61]
Let \(A\) be an involutive Banach algebra with a bounded approximate identity.
To any continuous positive linear functional \(\omega\), there corresponds
uniquely, within unitary equivalence, a representation
\(\{\pi_\omega,\mathcal H_\omega\}\) of \(A\) with a vector
\(\xi_\omega\) such that
\[
  [\pi_\omega(A)\xi_\omega]=\mathcal H_\omega
\]
and
\[
  \omega(x)=(\pi_\omega(x)\xi_\omega\mid\xi_\omega),
  \qquad x\in A.
\]
\end{theorem}

\begin{proof}
Uniqueness is obtained by sending \(\pi_\omega(x)\xi_\omega\) in one cyclic
triple to \(\pi'_\omega(x)\xi'_\omega\) in the other and using the displayed
coefficient identity to prove that this map is well-defined, isometric, and
intertwining.  Existence is obtained from the quotient by the left kernel of
\(\omega\), completed in the inner product
\((\eta_\omega(x)\mid\eta_\omega(y))=\omega(y^*x)\); left multiplication
gives bounded operators, and Riesz representation produces the vector
\(\xi_\omega\) implementing \(\omega\).
\end{proof}

\begin{definition}[GNSCyclicRepresentationTerminology, Takesaki I Definition I.9.15, PDF p.~61]
The representation \(\{\pi_\omega,\mathcal H_\omega\}\) constructed in
Theorem I.9.14 is called the cyclic representation of \(A\) induced by
\(\omega\).  Sometimes it is denoted by
\(\{\pi_\omega,\mathcal H_\omega,\xi_\omega\}\) to indicate the vector
corresponding to \(\omega\).  The construction is called the
Gelfand--Naimark--Segal construction.  In general, if a representation
\(\{\pi,\mathcal H\}\) of \(A\) admits a vector \(\xi\) such that
\([\pi(A)\xi]=\mathcal H\), then \(\{\pi,\mathcal H\}\) is said to be cyclic
and \(\xi\) is called a cyclic vector for \(\pi\).
\end{definition}

\begin{definition}[VonNeumannAlgebraGeneratedAndSpatialIso, Takesaki I Definition II.3.2, PDF p.~92]
A von Neumann algebra on \(\mathcal H\) is a \(*\)-subalgebra
\(\mathcal M\) of \(\mathcal L(\mathcal H)\) such that
\[
  \mathcal M=\mathcal M''.
\]
A factor is a von Neumann algebra \(\mathcal M\) with trivial center,
\(\mathcal M\cap\mathcal M'=\mathbb C1\).  If a subset \(Y\) of
\(\mathcal L(\mathcal H)\) is invariant under the \(*\)-operation, then
\(Y''\) is the smallest von Neumann algebra containing \(Y\), and is called
the von Neumann algebra generated by \(Y\).  If
\(\{\mathcal M_1,\mathcal H_1\}\) and
\(\{\mathcal M_2,\mathcal H_2\}\) are von Neumann algebras and there exists
an isometry \(U\) of \(\mathcal H_1\) onto \(\mathcal H_2\) such that
\[
  U\mathcal M_1U^*=\mathcal M_2,
\]
then they are said to be spatially isomorphic.
\end{definition}

\begin{theorem}[DoubleCommutationTheorem, Takesaki I Theorem II.3.9, PDF pp.~94--95]
Suppose \(\mathcal M\) is a \(*\)-algebra of operators on a Hilbert space
\(\mathcal H\).  If \(\mathcal M\) is \(\sigma\)-strongly* closed, then:
\[
\begin{array}{ll}
\textnormal{(i)}&
\textnormal{there exists a greatest projection }e\in\mathcal M,\textnormal{
the projection onto }[\mathcal M\mathcal H],\\
&
x=xe=ex\textnormal{ for every }x\in\mathcal M;\\[2mm]
\textnormal{(ii)}&
\mathcal M''\textnormal{ is the set of all operators }x+\alpha 1,
\ x\in\mathcal M,\ \alpha\in\mathbb C.
\end{array}
\]
Therefore, if \(\mathcal M\) is nondegenerate, then \(\mathcal M\) is a von
Neumann algebra, that is, \(\mathcal M=\mathcal M''\).
\end{theorem}

\begin{proof}
In the nondegenerate case, amplify to a countable direct sum and use the
approximation lemma for elements of \(\mathcal M''\) to obtain
\(\sigma\)-strong approximation by elements of \(\mathcal M\).  Since
\(\mathcal M\) is \(\sigma\)-strongly* closed, this gives
\(\mathcal M=\mathcal M''\).  In the degenerate case, the \(\sigma\)-weak
compactness of the unit ball gives an identity projection \(e\), and the
orthogonal complement contributes only scalar multiples of the identity.
\end{proof}

\begin{definition}[InfiniteTensorProductVNStatesAndVectors, Takesaki III Definition XIV.1.6, PDF pp.~106--107]
For von Neumann algebras \(M_n\) with normal states \(\omega_n\), the product
state on the \(C^*\)-algebraic infinite tensor product is defined on elementary
eventually unital tensors by
\[
  \omega(x_1\otimes x_2\otimes\cdots\otimes x_n\otimes 1\otimes 1\otimes
  \cdots)
  =
  \omega_1(x_1)\omega_2(x_2)\cdots\omega_n(x_n).
\]
The von Neumann algebra generated by the cyclic representation of this product
state is called the infinite tensor product of \(\{M_n\}\) relative to
\(\{\omega_n\}\).  If Hilbert spaces \(\mathcal H_n\) and unit reference
vectors \(\xi_n\in\mathcal H_n\) are chosen, the inductive limit of the finite
tensor-product Hilbert spaces under
\(\eta\mapsto\eta\otimes\xi_{n+1}\) is completed to the infinite tensor
product Hilbert space relative to the reference vectors.
\end{definition}

\subsection*{Conventions}

\begin{convention}[State and representation convention]
All \(C^*\)-algebras in this note are unital unless explicitly stated
otherwise.  A state is a positive linear functional of norm one.  For a state
\(\omega\) on a \(C^*\)-algebra \(A\), write
\[
  (\cH_\omega,\pi_\omega,\OmegaVec_\omega)
\]
for a cyclic GNS triple.
\source{\path{blueprints/008-stabilizer-local-state-foundations.tex}, source
facts for Takesaki I, Definition I.9.4, Lemma I.9.9, and Proposition I.9.12;
Takesaki I, Theorem I.9.14 and Definition I.9.15, PDF pp.~59--61.}
\status{Convention for all targets in this note.}
\end{convention}

\subsection*{Hilbert Interfaces}

\begin{definition}[HilbertInterfaces, ID: HILB-001, Takesaki I Theorem I.9.14, PDF pp.~59--61, Takesaki III Definition XIV.1.6, PDF pp.~106--107]
The Hilbert-space interfaces needed for the stabilizer type-classification
formalization are:
\begin{enumerate}
\item the GNS Hilbert space \(\cH_\omega\) of a quasi-local state
\(\omega\);
\item finite tensor products \(\bigotimes_{i\in F}\mathbb C^2\) for finite
Pauli blocks;
\item incomplete or reference-vector infinite tensor products used as sourced
abstractions for product-state GNS identifications;
\item a fixed separable infinite-dimensional Hilbert space, later usable for
model \(B(\ell^2)\)-type factors.
\end{enumerate}
The formalization may rely on mathlib for ordinary Hilbert spaces, bounded
operators, finite tensor products where available, and GNS carrier types; the
infinite tensor-product Hilbert-space facts remain narrowly sourced
abstractions until implemented.
\source{Takesaki I, Theorem I.9.14, PDF pp.~59--61, for GNS Hilbert spaces;
Takesaki III, Definition XIV.1.6 and surrounding construction, PDF
pp.~106--107, for infinite tensor product Hilbert spaces relative to reference
vectors; \path{agents/2026-06-21-mathlib-operator-algebra-coverage-audit.md}
for mathlib Hilbert-space coverage.}
\depends{\texttt{STATE-003}.}
\lean{Formalization.OperatorAlgebra.HilbertInterfaces}
\status{blueprint_note_complete}
\end{definition}

\subsection*{GNS Package and Transport}

\begin{theorem}[GNSPackage, ID: GNS-001, Takesaki I Theorem I.9.14 and Definition I.9.15, PDF pp.~59--61]
Let \(A\) be a unital \(C^*\)-algebra and let \(\omega\) be a state on
\(A\).  There is a Hilbert space \(\cH_\omega\), a unital
\(*\)-representation
\[
  \pi_\omega:A\to\cB(\cH_\omega),
\]
and a cyclic unit vector \(\OmegaVec_\omega\in\cH_\omega\) such that
\[
  \omega(a)=
  \langle \pi_\omega(a)\OmegaVec_\omega,\OmegaVec_\omega\rangle
\]
for every \(a\in A\).  The cyclic triple is unique up to unitary equivalence.
\source{Takesaki I, Section I.9, Theorem I.9.14, PDF pp.~59--61, and
Definition I.9.15, PDF p.~61.  Lean implementation target:
\path{PositiveLinearMap.GNS} and \path{PositiveLinearMap.gnsStarAlgHom} from
the mathlib coverage audit.}
\depends{\texttt{HILB-001}, \texttt{STATE-003}.}
\lean{Formalization.GNS.GNSPackage}
\status{blueprint_note_complete}
\end{theorem}

\begin{proof}
A state on a unital \(C^*\)-algebra is a continuous positive linear functional
by the state and continuity facts already recorded in
\path{blueprints/008-stabilizer-local-state-foundations.tex}.  The unital
algebra has a bounded approximate identity, namely \(1\).  Takesaki I,
Theorem I.9.14 therefore applies and gives the cyclic representation, cyclic
vector, reconstruction formula, and uniqueness up to unitary equivalence.
Definition I.9.15 supplies the GNS terminology.
\end{proof}

\begin{theorem}[CanonicalProjectionCodeGNS, ID: GNS-004, Takesaki I Theorem I.9.14 and Definition I.9.15, PDF pp.~59--61]
Let \(S=(\Gamma,W,\chi)\) be a stabilizer system with canonical finite-local
projection-code states from \(\texttt{STATE-002}\), and let
\[
  \omega_S:\cA_\Gamma\to\mathbb C
\]
be the induced quasi-local canonical state from \(\texttt{STATE-003}\).  The
canonical represented Hilbert-space data for this stabilizer system are defined
to be the GNS triple
\[
  (\cH_{\omega_S},\pi_{\omega_S},\OmegaVec_{\omega_S})
\]
of \(\omega_S\).  Thus \(\cH_{\omega_S}\) is a Hilbert space,
\[
  \pi_{\omega_S}:\cA_\Gamma\to\cB(\cH_{\omega_S})
\]
is a unital \(*\)-representation, \(\OmegaVec_{\omega_S}\) is a cyclic unit
vector, and for every \(a\in\cA_\Gamma\),
\[
  \omega_S(a)=
  \langle \pi_{\omega_S}(a)\OmegaVec_{\omega_S},
          \OmegaVec_{\omega_S}\rangle .
\]
For every region \(\Lambda\subseteq\Gamma\), the represented region von
Neumann algebra in this project is formed using this representation:
\[
  \cM_\Lambda(S)
  :=
  \pi_{\omega_S}(\cA_\Lambda)''\subseteq\cB(\cH_{\omega_S}).
\]
\source{The state \(\omega_S\) is supplied by
\path{blueprints/008-stabilizer-local-state-foundations.tex},
\texttt{STATE-003}.  Takesaki I, Theorem I.9.14 and Definition I.9.15, PDF
pp.~59--61, supply the GNS cyclic triple and reconstruction formula.}
\depends{\texttt{STATE-002}, \texttt{STATE-003}, \texttt{GNS-001}.}
\lean{Formalization.GNS.canonicalProjectionCodeGNS}
\status{blueprint_note_complete}
\end{theorem}

\begin{proof}
By \(\texttt{STATE-003}\), the compatible finite-local projection-code
functionals extend uniquely to a state \(\omega_S\) on the quasi-local
\(C^*\)-algebra \(\cA_\Gamma\).  Applying \(\texttt{GNS-001}\), equivalently
Takesaki I, Theorem I.9.14, to the pair \((\cA_\Gamma,\omega_S)\) gives a
Hilbert space, a unital cyclic representation, a cyclic unit vector, and the
displayed coefficient formula.  Definition I.9.15 identifies this cyclic
triple as the GNS representation induced by \(\omega_S\).  The final displayed
definition is therefore not an additional theorem: it records that all region
algebras below are represented on the specific GNS Hilbert space of the
canonical projection-code state.
\end{proof}

\begin{lemma}[GNSUnitaryOfStateIso, ID: GNS-002, Takesaki I Theorem I.9.14, PDF pp.~59--61]
Let \(\alpha:A\to B\) be a unital \(*\)-isomorphism of \(C^*\)-algebras, and
let \(\varphi\) be a state on \(B\).  The GNS representation of
\(\varphi\circ\alpha\) is unitarily equivalent to the pulled-back cyclic
representation
\[
  a\mapsto \pi_\varphi(\alpha(a))
\]
on \(\cH_\varphi\).  Under this unitary, the von Neumann algebra generated by
the represented image of any subalgebra \(A_0\subseteq A\) is spatially
isomorphic to the von Neumann algebra generated by the represented image of
\(\alpha(A_0)\subseteq B\).
\source{Takesaki I, Theorem I.9.14, PDF pp.~59--61, for uniqueness of cyclic
representations; Takesaki I, Definition II.3.2, PDF p.~92, for spatial
isomorphism and generated von Neumann algebras.}
\depends{\texttt{GNS-001}, \texttt{VN-001}, \texttt{SYMP-002}.}
\lean{Formalization.GNS.gnsUnitaryOfStateIso}
\status{blueprint_note_complete}
\end{lemma}

\begin{proof}
The triple
\[
  (\cH_\varphi,\pi_\varphi\circ\alpha,\OmegaVec_\varphi)
\]
is cyclic for \(A\) and implements the state \(\varphi\circ\alpha\), because
\[
  \langle \pi_\varphi(\alpha(a))\OmegaVec_\varphi,\OmegaVec_\varphi\rangle
  =\varphi(\alpha(a)).
\]
By GNS uniqueness, there is a unitary intertwiner from the canonical GNS
triple of \(\varphi\circ\alpha\) to this pulled-back triple.  If \(A_0\) is a
subalgebra, the intertwining identity sends the represented set
\(\pi_{\varphi\circ\alpha}(A_0)\) to \(\pi_\varphi(\alpha(A_0))\).  Unitarily
conjugating a set also conjugates its commutant and hence its double
commutant, so the generated von Neumann algebras are spatially isomorphic.
\end{proof}

\begin{theorem}[ProductStateGNS, ID: GNS-003, Takesaki III Definition XIV.1.6, PDF pp.~106--107]
For the product state obtained in \texttt{PROD-001}, the GNS representation is
the product-state cyclic representation built from the one-site GNS
representations.  For a countable product, the Hilbert space is the incomplete
tensor product relative to the one-site cyclic vectors, and the represented
von Neumann algebra is the infinite tensor product generated by the one-site
represented algebras relative to those reference states or vectors.  For a
finite product, this reduces to the ordinary finite tensor-product
representation.
\source{Takesaki III, Definition XIV.1.6 and surrounding paragraphs, PDF
pp.~106--107, for product states, cyclic representations, and infinite tensor
products relative to reference states/vectors; \texttt{PROD-001} in
\path{blueprints/009-internal-stabilizer-reduction.tex} for the stabilizer
product state.}
\depends{\texttt{HILB-001}, \texttt{GNS-001}, \texttt{GNS-002},
\texttt{PROD-001}.}
\lean{Formalization.GNS.productStateGNS}
\status{blueprint_note_complete}
\end{theorem}

\begin{proof}
The state from \texttt{PROD-001} is a tensor product of one-site vector states
on fixed modes and one-site normalized traces on defect modes.  Takesaki III's
construction forms the product state on the \(C^*\)-algebraic infinite tensor
product, then takes its cyclic representation and generated von Neumann
algebra.  The one-site cyclic vectors are exactly the reference vectors for
the incomplete tensor product Hilbert space.  Since the finite Pauli local
algebras form a cofinal local net, equality on all finite tensor factors
identifies this representation with the GNS representation of the product
state.  The finite case is the same argument without a limiting step.
\end{proof}

\subsection*{Represented Region Von Neumann Algebras}

\begin{definition}[ConcreteCommutant, ID: VN-004, Takesaki I Definition II.3.2, PDF p.~92]
Let \(\cH\) be a Hilbert space and let
\(S\subseteq\cB(\cH)\).  The commutant of \(S\) is
\[
  S'
  :=
  \{T\in\cB(\cH): Ts=sT\ \text{for every }s\in S\}.
\]
The double commutant of \(S\) is
\[
  S'' := (S')'.
\]
When \(S\) is invariant under adjoint, \(S'\) and \(S''\) are unital
\(*\)-subalgebras of \(\cB(\cH)\).
\source{Takesaki I, Definition II.3.2, PDF p.~92, for commutants, double
commutants, and von Neumann algebras defined by the double-commutant
condition.}
\depends{\texttt{HILB-001}.}
\lean{Formalization.OperatorAlgebra.commutant; Formalization.OperatorAlgebra.doubleCommutant}
\status{blueprint_note_complete}
\end{definition}

\begin{proof}[Informal well-formedness]
The identity operator commutes with every element of \(S\).  If two operators
commute with every \(s\in S\), then so do their linear combinations and
products, so \(S'\) is a unital subalgebra.  If \(S=S^*\), then from
\(Ts=sT\) one obtains \(sT^*=T^*s\) by applying adjoints to the equation for
\(s^*\in S\); hence \(S'\) is closed under adjoint.  Applying the same
definition once more gives \(S''\).
\end{proof}

\begin{definition}[GeneratedVonNeumannAlgebra, ID: VN-005, Takesaki I Definition II.3.2, PDF p.~92]
Let \(S\subseteq\cB(\cH)\) be invariant under adjoint.  The concrete von
Neumann algebra generated by \(S\) is
\[
  W^*(S) := S''.
\]
For a unital \(*\)-subalgebra \(A\leq\cB(\cH)\), this specializes to
\[
  W^*(A)=A''.
\]
\source{Takesaki I, Definition II.3.2, PDF p.~92: if \(Y\subseteq
\mathcal L(\mathcal H)\) is invariant under the \(*\)-operation, then
\(Y''\) is the smallest von Neumann algebra containing \(Y\), and is called
the von Neumann algebra generated by \(Y\).}
\depends{\texttt{VN-004}.}
\lean{Formalization.OperatorAlgebra.generatedVonNeumannAlgebra}
\status{blueprint_note_complete}
\end{definition}

\begin{proof}[Informal well-formedness]
The adjoint-invariance hypothesis is the one used in Takesaki's generated
von Neumann algebra statement.  It is satisfied for the image of any
\(*\)-subalgebra under a \(*\)-representation, in particular for
\(\pi_{\omega_S}(\cA_\Lambda)\).  The notation \(W^*(S)\) is only a
Lean-facing wrapper for the double commutant \(S''\).
\end{proof}

\begin{proposition}[GeneratedVonNeumannAlgebraMonotoneMinimal, ID: VN-006, Takesaki I Definition II.3.2, PDF p.~92]
Let \(S,T\subseteq\cB(\cH)\) be adjoint-invariant sets.
\begin{enumerate}
\item If \(S\subseteq T\), then \(W^*(S)\subseteq W^*(T)\).
\item If \(\cN\subseteq\cB(\cH)\) is a concrete von Neumann algebra and
  \(S\subseteq\cN\), then \(W^*(S)\subseteq\cN\).
\item Consequently \(W^*(S)=S''\) is the smallest concrete von Neumann
  algebra on \(\cH\) containing \(S\).
\end{enumerate}
\source{Takesaki I, Definition II.3.2, PDF p.~92, for the generated
von Neumann algebra as \(S''\) and for the defining equality
\(\cN=\cN''\).}
\depends{\texttt{VN-004}, \texttt{VN-005}.}
\lean{Formalization.OperatorAlgebra.generatedVN_mono; Formalization.OperatorAlgebra.generatedVN_le}
\status{blueprint_note_complete}
\end{proposition}

\begin{proof}
Commutants reverse inclusions: if \(S\subseteq T\) and an operator commutes
with every element of \(T\), then it commutes with every element of \(S\).
Thus \(T'\subseteq S'\), and applying the commutant operation again gives
\(S''\subseteq T''\), which is monotonicity of \(W^*(-)\).

If \(S\subseteq\cN\), then \(\cN'\subseteq S'\), hence
\[
  S''\subseteq \cN''=\cN
\]
because \(\cN\) is a von Neumann algebra.  Since every \(s\in S\) commutes
with every operator in \(S'\), one has \(S\subseteq S''=W^*(S)\).  Therefore
any concrete von Neumann algebra containing \(S\) also contains \(W^*(S)\),
which is precisely the minimality claim.
\end{proof}

\begin{theorem}[UnitalStarSubalgebraDoubleCommutantClosures, ID: VN-DCT-001, derived from Takesaki I Theorem II.2.6, Lemma II.3.8, and proof of Theorem II.3.9, PDF pp.~90, 94--95]
Let \(A\leq\cB(\cH)\) be a unital \(*\)-subalgebra.  Then the double
commutant, strong-operator closure, and weak-operator closure of \(A\) agree:
\[
  A''
  =
  \overline{A}^{\,\mathrm{SOT}}
  =
  \overline{A}^{\,\mathrm{WOT}}
  \subseteq \cB(\cH).
\]
Consequently, if a net or sequence of elements of \(A\) converges strongly to
\(x\), then \(x\in A''\).  In particular, a concrete von Neumann algebra
\(\cM=\cM''\) is closed under strong-operator and weak-operator limits of
nets from \(\cM\).
\source{This is a project-derived bicommutant theorem, not the verbatim
closed-algebra statement alone.  Takesaki I, Theorem II.3.9, PDF pp.~94--95,
states the double commutation theorem for \(\sigma\)-strongly* closed
\(*\)-algebras.  Its proof paragraph beginning ``At first, suppose
\(\mathcal M\) is nondegenerate'' on PDF p.~94 and continuing on PDF p.~95
uses Lemma II.3.8, PDF p.~94, to approximate an element of the double
commutant in the \(\sigma\)-strong topology by elements of the algebra.
Takesaki I, Theorem II.2.6, PDF p.~90, compares weak, strong,
\(\sigma\)-weak, and \(\sigma\)-strong closures for convex operator sets and
is the closure-comparison input used to pass from that approximation argument
to the displayed unital-subalgebra closure equality.}
\depends{\texttt{VN-004}, \texttt{VN-005}.}
\lean{Formalization.OperatorAlgebra.doubleCommutant_eq_sotClosure;
Formalization.OperatorAlgebra.doubleCommutant_eq_wotClosure}
\status{blueprint_note_complete}
\end{theorem}

\begin{proof}
Because \(A\) is unital, it is nondegenerate as an algebra of operators:
\([A\cH]=\cH\).  The set \(A''\) is weak-operator closed, since commutation
with a fixed bounded operator is expressed by weak-operator closed linear
equations.  It contains \(A\), so
\[
  \overline{A}^{\,\mathrm{WOT}}\subseteq A''.
\]
The strong-operator topology is finer than the weak-operator topology, hence
\[
  \overline{A}^{\,\mathrm{SOT}}\subseteq
  \overline{A}^{\,\mathrm{WOT}}.
\]

For the reverse inclusion, this target uses the bicommutant approximation
inside Takesaki's proof rather than treating the closed-algebra transcript as
the direct statement.  Let \(B\) be the \(\sigma\)-strong closure of \(A\).
By Theorem II.2.6, equivalently taking the \(\sigma\)-strong* closure gives
the same closed convex \(*\)-algebra.  Since \(A\) is unital, \(B\) is
nondegenerate.  The commutant is unchanged by this closure:
if an operator commutes with every element of \(A\), then it commutes with
every \(\sigma\)-strong limit of elements of \(A\), so \(A'=B'\) and hence
\(A''=B''\).

Now apply the nondegenerate part of Takesaki I, Theorem II.3.9 to \(B\).  The
proof on PDF pp.~94--95 uses Lemma II.3.8 and the amplification argument to
show that every element of \(B''\) is \(\sigma\)-strongly approximated by
elements of \(B\); because \(B\) is \(\sigma\)-strongly* closed, this yields
\(B=B''\).  Thus \(A''=B\).  Theorem II.2.6 then identifies this same closed
linear subspace with the strong-operator and weak-operator closures of \(A\).
Therefore
\[
  A''\subseteq\overline{A}^{\,\mathrm{SOT}}
  \subseteq\overline{A}^{\,\mathrm{WOT}}
  \subseteq A'',
\]
which proves the displayed equality.  The strong-limit membership statement
is immediate from \(A''=\overline{A}^{\,\mathrm{SOT}}\).  Taking
\(A=\cM\) and using \(\cM=\cM''\) gives the closure of concrete von Neumann
algebras under strong and weak limits from themselves.
\end{proof}

\begin{definition}[RegionVNPackage, ID: VN-001, Takesaki I Definition II.3.2, PDF p.~92]
Let \(\omega\) be a state on the quasi-local algebra \(\cA_\Gamma\), with GNS
representation \((\cH_\omega,\pi_\omega,\OmegaVec_\omega)\).  For a region
\(\Lambda\subseteq\Gamma\), let \(\cA_\Lambda\subseteq\cA_\Gamma\) be the
region \(C^*\)-subalgebra.  The represented region von Neumann algebra is
\[
  \cM_\Lambda(\omega)
  :=
  \pi_\omega(\cA_\Lambda)''\subseteq\cB(\cH_\omega).
\]
Equivalently, it is the smallest concrete von Neumann algebra on
\(\cH_\omega\) containing the represented region algebra
\(\pi_\omega(\cA_\Lambda)\).  For the canonical stabilizer system \(S\), this
definition uses the GNS triple
\((\cH_{\omega_S},\pi_{\omega_S},\OmegaVec_{\omega_S})\) fixed in
\(\texttt{GNS-004}\), and we write
\[
  \cM_\Lambda(S):=
  \pi_{\omega_S}(\cA_\Lambda)''\subseteq\cB(\cH_{\omega_S}).
\]
\source{Takesaki I, Definition II.3.2, PDF p.~92, for
von Neumann algebras as double commutants and \(Y''\) as the generated von
Neumann algebra.}
\depends{\texttt{GNS-001}, \texttt{GNS-004}, \texttt{QL-002}, \texttt{VN-005}.}
\lean{Formalization.OperatorAlgebra.RegionVNPackage}
\status{blueprint_note_complete}
\end{definition}

\begin{proof}[Informal well-formedness]
The region algebra \(\cA_\Lambda\) is a \(C^*\)-subalgebra by
\(\texttt{QL-002}\).  Since \(\pi_\omega\) is a \(*\)-representation, its
image \(\pi_\omega(\cA_\Lambda)\) is invariant under adjoint and contains the
identity of \(\cB(\cH_\omega)\).  Therefore the generated von Neumann algebra
\(W^*(\pi_\omega(\cA_\Lambda))=\pi_\omega(\cA_\Lambda)''\) is well-defined
by \(\texttt{VN-005}\).  In the canonical case, \(\texttt{GNS-004}\) supplies
the specific state and representation used in this formula.
\end{proof}

\begin{lemma}[RepresentedRegionVNIsotony, ID: VN-007]
Let \(\omega\) be a state on \(\cA_\Gamma\), and let
\(\Lambda_1\subseteq\Lambda_2\subseteq\Gamma\).  Then
\[
  \cM_{\Lambda_1}(\omega)\subseteq \cM_{\Lambda_2}(\omega).
\]
In particular, for the canonical projection-code state \(\omega_S\),
\[
  \cM_{\Lambda_1}(S)\subseteq \cM_{\Lambda_2}(S).
\]
\source{Project-local consequence of the region-algebra isotony in
\(\texttt{QL-002}\) and generated-von-Neumann-algebra monotonicity in
\(\texttt{VN-006}\).}
\depends{\texttt{QL-002}, \texttt{GNS-004}, \texttt{VN-001}, \texttt{VN-006}.}
\lean{Formalization.OperatorAlgebra.regionVN_isotone}
\status{blueprint_note_complete}
\end{lemma}

\begin{proof}
Region isotony gives \(\cA_{\Lambda_1}\subseteq\cA_{\Lambda_2}\).  Applying
the representation \(\pi_\omega\) preserves inclusion:
\[
  \pi_\omega(\cA_{\Lambda_1})
  \subseteq
  \pi_\omega(\cA_{\Lambda_2}).
\]
Both represented sets are adjoint-invariant.  By the monotonicity of
generated von Neumann algebras from \(\texttt{VN-006}\), their double
commutants satisfy
\[
  \pi_\omega(\cA_{\Lambda_1})''
  \subseteq
  \pi_\omega(\cA_{\Lambda_2})''.
\]
This is the asserted isotony, and substituting \(\omega=\omega_S\) gives the
canonical projection-code case.
\end{proof}

\begin{lemma}[GeneratedVNUnitaryConj, ID: VN-008, Takesaki I Definition II.3.2, PDF p.~92]
Let \(U:\cH_1\to\cH_2\) be a unitary and let
\(S\subseteq\cB(\cH_1)\) be adjoint-invariant.  Put
\[
  USU^*=\{UsU^*:s\in S\}\subseteq\cB(\cH_2).
\]
Then
\[
  (USU^*)''=US''U^*.
\]
Equivalently,
\[
  W^*(USU^*)=U\,W^*(S)\,U^*.
\]
Thus unitary conjugacy is a spatial isomorphism between the generated concrete
von Neumann algebras.
\source{Takesaki I, Definition II.3.2, PDF p.~92, for spatial isomorphism
and for generated von Neumann algebras as double commutants.}
\depends{\texttt{VN-004}, \texttt{VN-005}.}
\lean{Formalization.OperatorAlgebra.generatedVN_unitaryConj}
\status{blueprint_note_complete}
\end{lemma}

\begin{proof}
For \(T\in\cB(\cH_2)\), the equation \(T(UsU^*)=(UsU^*)T\) for all
\(s\in S\) is equivalent, after multiplying by \(U^*\) on the left and \(U\)
on the right, to
\[
  (U^*TU)s=s(U^*TU)
  \qquad\text{for all }s\in S.
\]
Therefore
\[
  (USU^*)'=US'U^*.
\]
Applying the same identity to the adjoint-invariant set \(S'\) gives
\[
  (USU^*)''=(US'U^*)'=US''U^*.
\]
By the definition \(W^*(S)=S''\), this is the displayed generated-algebra
identity.  Takesaki's spatial-isomorphism definition is exactly the existence
of such a unitary implementing \(U W^*(S) U^*=W^*(USU^*)\).
\end{proof}

\begin{lemma}[RegionVNConj, ID: VN-003, Takesaki I Definition II.3.2 and Theorem II.3.9, PDF pp.~92--95]
Let \(U:\cH_1\to\cH_2\) be a unitary intertwining two represented subalgebras:
\[
  U\pi_1(a)U^*=\pi_2(\alpha(a)),\qquad a\in A_0.
\]
Then
\[
  U\,\pi_1(A_0)''\,U^*=\pi_2(\alpha(A_0))''.
\]
In particular, represented region von Neumann algebras related by the GNS
transport unitary of \texttt{GNS-002} are spatially isomorphic.
\source{Takesaki I, Definition II.3.2, PDF p.~92, for commutants, generated
von Neumann algebras, and spatial isomorphism; Theorem II.3.9, PDF pp.~94--95,
for the double-commutation theorem background.}
\depends{\texttt{VN-001}, \texttt{VN-008}, \texttt{GNS-002}.}
\lean{Formalization.OperatorAlgebra.regionVN_conj}
\status{blueprint_note_complete}
\end{lemma}

\begin{proof}
For any set \(S\subseteq\cB(\cH_1)\), unitary conjugation satisfies
\[
  (USU^*)'=US'U^*.
\]
Indeed, an operator on \(\cH_2\) commutes with every \(UsU^*\) exactly when
its conjugate by \(U^*\) commutes with every \(s\in S\).  Applying the same
identity once more gives
\[
  (USU^*)''=US''U^*.
\]
Taking \(S=\pi_1(A_0)\) and using the intertwining hypothesis gives the
displayed equality.  Takesaki's Definition II.3.2 identifies these double
commutants as the generated von Neumann algebras and calls the unitary
conjugacy a spatial isomorphism.
\end{proof}

\begin{theorem}[DecreasingProjectionStrongLimit, ID: PROJ-LIM-001, Kadison--Ringrose I Proposition 2.5.6 and Corollary 2.5.7, local PDF pp.~121--124; Takesaki I Proposition V.1.1, PDF p.~310]
Let \((p_n)_{n\in\mathbb N}\) be a decreasing sequence of orthogonal
projections on a Hilbert space \(\cH\), meaning
\[
  p_n=p_n^*=p_n^2,
  \qquad
  p_{n+1}\leq p_n
  \quad\text{for all }n.
\]
Let
\[
  K=\bigcap_{n\in\mathbb N}\operatorname{Ran}(p_n),
\]
and let \(p\) be the orthogonal projection onto the closed subspace \(K\).
Then:
\begin{enumerate}
\item \(p_n\to p\) in the strong operator topology;
\item \(\operatorname{Ran}(p)=\bigcap_n\operatorname{Ran}(p_n)\);
\item \(p\leq p_n\) for every \(n\);
\item \(p=\bigwedge_n p_n\), the infimum in the projection order.
\end{enumerate}
If \(\cM\subseteq\cB(\cH)\) is a concrete von Neumann algebra and
\(p_n\in\cM\) for every \(n\), then \(p\in\cM\).
\source{Kadison--Ringrose, \emph{Fundamentals of the Theory of Operator
Algebras I}, Proposition 2.5.6 and Corollary 2.5.7, local PDF pp.~121--124
(printed pp.~112--115), for strong convergence of monotone projection nets to
their supremum/infimum.  Takesaki I, Proposition V.1.1, PDF p.~310, for the
complete projection lattice of a von Neumann algebra and the range-intersection
description of infima.  Membership in \(\cM\) uses
\(\texttt{VN-DCT-001}\).}
\depends{\texttt{VN-DCT-001}.}
\lean{Formalization.OperatorAlgebra.decreasingProjection_strongLimit}
\status{blueprint_note_complete}
\end{theorem}

\begin{proof}
The intersection \(K=\bigcap_n\operatorname{Ran}(p_n)\) is closed, because it
is an intersection of closed subspaces, so it has an orthogonal projection
\(p\).  Since \(K\subseteq\operatorname{Ran}(p_n)\) for every \(n\), the
projection order gives \(p\leq p_n\).

If \(q\) is any projection with \(q\leq p_n\) for all \(n\), then
\(\operatorname{Ran}(q)\subseteq\operatorname{Ran}(p_n)\) for all \(n\).
Hence \(\operatorname{Ran}(q)\subseteq K=\operatorname{Ran}(p)\), so
\(q\leq p\).  Thus \(p\) is the greatest lower bound
\(\bigwedge_n p_n\), and its range is exactly the displayed intersection.

Kadison--Ringrose I, Corollary 2.5.7 applies to the decreasing sequence,
viewed as a decreasing net indexed by \(\mathbb N\), and states that the
infimum projection \(p=\bigwedge_n p_n\) satisfies \(p_n x\to px\) in norm
for every \(x\in\cH\).  This is exactly strong-operator convergence
\(p_n\to p\).

Finally suppose \(p_n\in\cM\) for all \(n\) and \(\cM=\cM''\).  By
\(\texttt{VN-DCT-001}\), \(\cM\) is strong-operator closed.  Since the
sequence \(p_n\) converges strongly to \(p\), the limit projection \(p\)
belongs to \(\cM\).
\end{proof}

\subsection*{Lean Extraction Notes}

\begin{longtable}{@{}p{0.15\linewidth}p{0.35\linewidth}p{0.40\linewidth}@{}}
\toprule
Target & Intended file & Notes \\
\midrule
\endhead
\texttt{HILB-001} &
\texttt{Formalization/OperatorAlgebra/HilbertInterfaces.lean} &
Use mathlib Hilbert spaces and bounded operators; keep infinite tensor-product
Hilbert facts as sourced abstractions. \\
\texttt{GNS-001} &
\texttt{Formalization/GNS/Package.lean} &
Wrap \texttt{PositiveLinearMap.GNS} and
\texttt{PositiveLinearMap.gnsStarAlgHom}; expose cyclic-vector and
reconstruction fields only as needed. \\
\texttt{GNS-004} &
\texttt{Formalization/GNS/CanonicalProjectionCode.lean} &
Specialize the GNS package to the \texttt{STATE-003} canonical projection-code
state and make this the representation used by represented region algebras. \\
\texttt{GNS-002} &
\texttt{Formalization/GNS/Equivalence.lean} &
Project-local wrapper from GNS uniqueness under state transport by a
\(*\)-isomorphism. \\
\texttt{GNS-003} &
\texttt{Formalization/GNS/ProductState.lean} &
Likely starts as a sourced abstraction for countable infinite products; finite
blocks can use existing tensor-product infrastructure. \\
\texttt{VN-001} &
\texttt{Formalization/OperatorAlgebra/RegionVN.lean} &
Concrete \texttt{VonNeumannAlgebra H} generated by the represented region
image. \\
\texttt{VN-004} &
\texttt{Formalization/OperatorAlgebra/Commutant.lean} &
Direct concrete definitions of \(S'\) and \(S''\) for sets of bounded
operators. \\
\texttt{VN-005} &
\texttt{Formalization/OperatorAlgebra/GeneratedVN.lean} &
Define \(W^*(S)\) as \(S''\) for adjoint-invariant represented sets. \\
\texttt{VN-006} &
\texttt{Formalization/OperatorAlgebra/GeneratedVN.lean} &
Monotonicity and minimality of generated von Neumann algebras. \\
\texttt{VN-DCT-001} &
\texttt{Formalization/OperatorAlgebra/DoubleCommutant.lean} &
Lean-facing bicommutant theorem:
\(A''=\overline A^{\mathrm{SOT}}=\overline A^{\mathrm{WOT}}\) for unital
\(*\)-subalgebras of bounded operators. \\
\texttt{VN-007} &
\texttt{Formalization/OperatorAlgebra/RegionVN.lean} &
Represented-region isotony from region-algebra isotony and generated-VN
monotonicity. \\
\texttt{VN-008} &
\texttt{Formalization/OperatorAlgebra/GeneratedVN.lean} &
Unitary conjugacy for generated von Neumann algebras. \\
\texttt{VN-003} &
\texttt{Formalization/OperatorAlgebra/RegionVNEquivalence.lean} &
Commutant manipulation under unitary conjugation; should be mathlib-backed
once the represented sets are packaged correctly. \\
\texttt{PROJ-LIM-001} &
\texttt{Formalization/OperatorAlgebra/ProjectionLimits.lean} &
Decreasing projection strong limits, range intersection, infimum, and
membership in a von Neumann algebra by strong closure. \\
\bottomrule
\end{longtable}

\begin{thebibliography}{9}
\bibitem{TakesakiI}
M. Takesaki, \emph{Theory of Operator Algebras I}.  Local source:
\path{../sources/takesaki_1.pdf}.

\bibitem{TakesakiIII}
M. Takesaki, \emph{Theory of Operator Algebras III}.  Local source:
\path{../sources/Takesaki_III.pdf}.

\bibitem{KadisonRingroseI}
R. V. Kadison and J. R. Ringrose, \emph{Fundamentals of the Theory of
Operator Algebras I}.  Local source: \path{../sources/kadison_ringrose_1.pdf}.
\end{thebibliography}

\end{document}
